\chapter{Conclusion}
This study investigates the impact of a localized shock, Hurricane Ian, on remittance flows from U.S. states to Mexican municipalities. Our baseline analysis reveals that a 1\% increase in Florida exposure is associated with a 0.16\% immediate decline in remittance inflows. However, potential compensatory effects may confound these results. While exposure to California yields no significant spillover effects, exposure to Texas exhibits a significant amplification effect. Indeed, a 1 percentage-point increase in Texas exposure worsens the negative impact of the Florida shock. Specifically, municipalities highly exposed to Texas experience a larger decline in remittances in the quarters following the shock, whereas municipalities with low exposure show a small increase. One plausible explanation is that a strong Texas network enables easier out-migration from Florida. By relocating, these migrants forfeit the higher wages of Florida's post-hurricane reconstruction boom, leading to lower baseline incomes and thereby reducing their remittance capability. 
\par
Furthermore, our findings demonstrate that the effects on remittance flows from macro-level shocks, such as the Great Recession \parencite{caballero_2023}, cannot directly be generalized to localized events. Localized shocks like Hurricane Ian disrupt only a specific subset of the population and have a limited impact on the broader U.S. economy. This allows for compensatory behavior, as unaffected migrants across the country can alter their remittance patterns to compensate the drop in Florida remittances during the quarter of the shock or use their networks to absorb displaced workers. Therefore, while a localized shock provides a cleaner empirical setting to isolate direct effects, the resulting impact on overall remittance patterns is complex and closer to statistical insignificance. These network dynamics could be further explored by studying other localized economic or climate shocks. Finally, future research could substantially improve upon these results by utilizing bilateral state-to-municipality remittance data, which would allow researchers to explicitly separate the migration channels directly impacted by a shock from those that remain unaffected.
\par

